\label{ch:gravity}

\section{Chaîne constitutive du secteur gravitationnel : échelle chronologique \(108\), réduction tensorielle, holonomie et évaluation métrologique de \(G\)}

La gravité s'organise ici en quatre strates qu'il ne faut pas fusionner :

\begin{equation}
\label{eq:gravity-causal-chain}
\begin{aligned}
\text{structure chronologique CT108 et droites d'action}
&\longrightarrow \text{sélecteur de périodicité }\theta_{108}^{\rm circ}
\longrightarrow G_{108},\\
\text{incidence icosaédrique}
&\longrightarrow V_{12}\xrightarrow{P_5}V_5
\xrightarrow{\Gamma_5}\operatorname{STF}_3
\xrightarrow{\Lambda(n)}\mathcal H_{\rm TT}(n),\\
\text{route holonomique enrichie}
&\dashrightarrow G_{\rm hol},\\
\text{carte métrologique}
&\longrightarrow 6.67430\times10^{-11}\,
\mathrm{m^3\,kg^{-1}\,s^{-2}}.
\end{aligned}
\end{equation}

La première ligne exprime une condition intrinsèque et exacte de
périodicité. La deuxième construit
le support et sa réduction transverse sans trace. La troisième conserve une
frontière de construction explicite. La quatrième est une évaluation dans une
carte décimale déclarée : elle identifie la grandeur, mais ne sélectionne pas le caractère
gravitationnel.

\section{Droite de grandeur gravitationnelle et échelle chronologique \(108\)}

Soient \(c_{\rm int}>0\), \(t_0>0\),
\(\ell_0=c_{\rm int}t_0\) et une section positive
\(\hbar_a\) de la droite d'action, avec
\(a\in\{\mathrm{pre},\mathrm{ret}\}\). Le transducteur dimensionnel est

\begin{equation}
\label{eq:gravity-transducer}
G_a(\theta)
=\theta\frac{c_{\rm int}^{5}t_0^2}{\hbar_a}
=\theta\frac{c_{\rm int}^{3}\ell_0^2}{\hbar_a},
\qquad \theta>0.
\end{equation}

L'égalité des deux expressions utilise seulement
\(\ell_0=c_{\rm int}t_0\), et
\([G]=L^3M^{-1}T^{-2}\). La dimension est fixée par la droite ; le
problème mathématique est de sélectionner le caractère adimensionnel \(\theta\).

Le calendrier CT108 définit

\begin{equation}
\label{eq:gravity-clock108}
T_{108}=108t_0,
\qquad
\omega_{108}=\frac{2\pi}{108t_0},
\qquad
R_{108}=\frac{c_{\rm int}}{\omega_{108}}
=\frac{54}{\pi}\ell_0.
\end{equation}

Pour chaque feuillet d'action, soient

\begin{equation}
\label{eq:gravity-clock-energy-mass}
E_{108,a}=\hbar_a\omega_{108},
\qquad
m_{108,a}=\frac{E_{108,a}}{c_{\rm int}^2}.
\end{equation}

Alors, la longueur de Compton réduite coïncide avec le rayon du cycle :

\begin{equation}
\label{eq:gravity-compton-clock}
\bar\lambda_{C,a}
=\frac{\hbar_a}{m_{108,a}c_{\rm int}}
=R_{108}.
\end{equation}

\begin{theorem}[Unicité du sélecteur sous la condition de périodicité HMT]
\label{thm:gravity-circular-selector}
Avec la convention réduite
\(r_g=Gm/c_{\rm int}^2\) et la prémisse de périodicité HMT qui identifie le rayon
gravitationnel du quantum CT108 à son rayon de phase, la condition de clôture
\begin{equation}
\label{eq:gravity-closing-condition}
r_{g,a}(\theta)=R_{108}
\end{equation}
sélectionne une unique solution positive,
\begin{equation}
\label{eq:gravity-theta108}
\boxed{
\theta_{108}^{\rm circ}=\left(\frac{54}{\pi}\right)^2
=295.4525715010569415303525\ldots .}
\end{equation}
Dans cette réalisation,
\begin{equation}
\label{eq:gravity-four-lengths}
R_{108}=\bar\lambda_{C,a}=r_{g,a}=\ell_P^{\rm circ}.
\end{equation}
\end{theorem}

\begin{proof}
Par \cref{eq:gravity-transducer,eq:gravity-clock-energy-mass},
\[
r_{g,a}(\theta)
=\frac{G_a(\theta)m_{108,a}}{c_{\rm int}^2}
=\theta c_{\rm int}t_0^2\omega_{108}
=\theta\frac{\pi}{54}\ell_0.
\]
En le comparant à
\(R_{108}=(54/\pi)\ell_0\), on obtient nécessairement
\(\theta=(54/\pi)^2\). L'équation est linéaire en \(\theta\) et tous les
facteurs sont positifs, donc la solution est unique.
\end{proof}

L'égalité \(r_g=R_{108}\) ne se déduit pas de la dimensionnalité de
\cref{eq:gravity-transducer}. C'est la condition géométrique explicite de cette
réalisation HMT : action, Compton et retour périodique doivent déterminer une même
longueur. Le théorème démontre que, cette condition structurelle étant acceptée, le
caractère \(\theta\) est fixé sans utiliser la valeur SI de \(G\).

\begin{remark}[Convention de rayon]
Le théorème utilise le rayon gravitationnel réduit \(Gm/c^2\). Si on le remplace
par le rayon de Schwarzschild \(2Gm/c^2\), le sélecteur change d'un facteur
deux. Les deux expressions correspondent à des conditions géométriques distinctes,
c'est pourquoi la convention doit être déclarée avant le calcul.
\end{remark}

\section{Réciprocité \(G\leftrightarrow\hbar\) et conservation de l'aire de Planck}

\begin{corollary}[Famille de Planck associée à la périodicité temporelle de \(108\) phases]
\label{cor:gravity-planck-family}
Pour \(G_a=G_a(\theta_{108}^{\rm circ})\),
\begin{align}
\ell_P&=\frac{54}{\pi}\ell_0,
&t_P&=\frac{54}{\pi}t_0,
\label{eq:gravity-planck-length-time}\\
E_{P,a}&=\frac{\hbar_a}{t_P}
=\frac{\pi\hbar_a}{54t_0}
=\frac{h_a}{108t_0},
&\ell_P^2&=\frac{G_a\hbar_a}{c_{\rm int}^3}
=\theta_{108}^{\rm circ}\ell_0^2.
\label{eq:gravity-planck-energy-area}
\end{align}
Adem\'as,
\begin{equation}
\label{eq:gravity-planck-force-power-density}
F_{P,a}=\frac{c_{\rm int}^4}{G_a},
\qquad
P_{P,a}=\frac{c_{\rm int}^5}{G_a},
\qquad
\rho_{P,a}=\frac{\hbar_a}
{(\theta_{108}^{\rm circ})^2c_{\rm int}^5t_0^4}.
\end{equation}
\end{corollary}

\begin{proof}
Les trois premières identités découlent de
\(\ell_P=R_{108}\), \(t_P=\ell_P/c_{\rm int}\) et
\(h_a=2\pi\hbar_a\). Pour l'aire,
\[
\frac{G_a\hbar_a}{c_{\rm int}^3}
=\theta_{108}^{\rm circ}c_{\rm int}^2t_0^2
=\theta_{108}^{\rm circ}\ell_0^2.
\]
Les expressions restantes sont les compositions dimensionnelles de force,
de puissance et de masse par volume avec la même section.
\end{proof}

Soit
\begin{equation}
\label{eq:gravity-Ract}
R_{\rm act}:=\frac{\hbar_{\rm pre}}{\hbar_{\rm ret}}.
\end{equation}
Comme l'action apparaît au numérateur de la masse chronologique et au
dénominateur de \(G\),
\begin{equation}
\label{eq:gravity-twoway-covariance}
\frac{m_{108,\rm pre}}{m_{108,\rm ret}}=R_{\rm act},
\qquad
\frac{G_{\rm pre}}{G_{\rm ret}}=R_{\rm act}^{-1},
\qquad
G_a\hbar_a
=\theta_{108}^{\rm circ}c_{\rm int}^5t_0^2.
\end{equation}
Par conséquent, \(\ell_P^2=G_a\hbar_a/c^3\) ne change pas de feuillet. La gravité et
l'action s'orientent réciproquement, tandis que l'aire conserve la
géométrie commune. Cette annulation est l'interface de confluence vers l'opérateur
d'aire du \cref{ch:modular}.
